\documentclass[12pt, letterpaper]{article}
\usepackage[margin=1in]{geometry}
\usepackage{amsmath}
\usepackage{amssymb}
\usepackage{mathtools}
\usepackage{comment}
\usepackage{enumerate}
\usepackage{changepage}
\usepackage{amsthm}
\usepackage{bbm}
\usepackage{tikz-cd}
\usepackage{xcolor}
\usepackage{hyperref}

\setlength{\parindent}{0pt}

%% code from mathabx.sty and mathabx.dcl
\DeclareFontFamily{U}{mathx}{\hyphenchar\font45}
\DeclareFontShape{U}{mathx}{m}{n}{
      <5> <6> <7> <8> <9> <10>
      <10.95> <12> <14.4> <17.28> <20.74> <24.88>
      mathx10
      }{}
\DeclareSymbolFont{mathx}{U}{mathx}{m}{n}
\DeclareFontSubstitution{U}{mathx}{m}{n}
\DeclareMathAccent{\widecheck}{0}{mathx}{"71}
\DeclareMathAccent{\wideparen}{0}{mathx}{"75}

\def\cs#1{\texttt{\char`\\#1}}

\allowdisplaybreaks

\DeclareMathOperator{\Id}{Id}
\DeclareMathOperator{\im}{im}
\renewcommand{\Re}{\text{Re}}
\renewcommand{\Im}{\text{Im}}
\newcommand{\D}{\mathbf{D}}
\newcommand{\0}{\mathbf{0}}
\newcommand{\1}{\mathbf{1}}
\newcommand{\loc}{\text{loc}}
\DeclareMathOperator{\dist}{dist}
\DeclareMathOperator{\Span}{Span}
\DeclareMathOperator{\sign}{sign}
\DeclareMathOperator{\supp}{supp}
\DeclareMathOperator{\op}{op}
\DeclareMathOperator{\tr}{tr}
\DeclareMathOperator{\x}{\mathbf{x}}
\DeclareMathOperator{\y}{\mathbf{y}}
\DeclareMathOperator{\z}{\mathbf{z}}

\title{Math 104 Homework 3}
\author{Due Date: Friday, July 17 at 11:59pm}
\date{}

\begin{document}

\maketitle

\textbf{Instructions:} Submit pdf solutions in LaTeX (preferred), or \textit{neatly} handwritten. Unless otherwise stated, you may cite any result from class or from Rudin's book, but all other assertions you make must be proved rigorously. At the top of your submission, please write down the names of the other students you collaborated with and any outside sources you consulted.

\vspace{0.5cm}

\begin{enumerate}
    \item 
    Let $(X, d)$ be a metric space. A subset $E \subset X$ is called \textit{dense} in $E$ if, for any open $O \subset X$, we have $O \cap E \neq \varnothing$. A metric space $(X, d)$ is called \textit{separable} if it has a countable dense subset.
    \begin{enumerate}
        \item 
        Show that $\mathbb{Q}$ is dense in $\mathbb{R}$, and conclude that $\mathbb{R}$ is separable.

        \item
        Show that a subset $E$ is dense in $X$ if and only if $\overline{E} = X$.

        \item 
        By following these steps, prove that any compact metric space $(X, d)$ is separable:
        \begin{enumerate}
            \item 
            Using problem 5(c) on homework 2, show that for each $n \in \mathbb{N}$, the following is true: there is a finite set $S_n \subset X$ so that every point of $X$ is within distance $\frac{1}{n}$ of some point of $S_n$.

            \item 
            Let $S = \bigcup\limits_{n = 1}^{\infty} S_n$. Show that $S$ is countable.

            \item 
            Show that every point of $X$ is either in $S$ or a limit point of $S$. Conclude that $S$ is dense in $X$, and hence that $X$ is separable.
            
        \end{enumerate}
        
    \end{enumerate}

    \vspace{0.5cm}

    \item 
    Using \textbf{only} the definitions of $\lim$, $\limsup$, and $\liminf$, prove the following statements.
    \begin{enumerate}
        \item 
        $\displaystyle{\lim_{n \to \infty} \left[\sqrt{n + 1} - \sqrt{n}\right] = 0}.$ \quad \textit{Hint:} Use difference of squares.\\

        \item 
        $\displaystyle{\limsup_{n \to \infty} \; (-2)^n = \infty}$.\\

        \item 
        $\displaystyle{\limsup_{n \to \infty} \; \left[-\sqrt{n}\right]} = -\infty.$\\
        

        \item 
        $\displaystyle{\liminf_{n \to \infty} \left[\frac{1}{n} + (1 + (-1)^n)n\right] = 0}.\\$
        
    \end{enumerate}

    \newpage

    \item 
    Let $\{s_n\}$ be an arbitrary sequence in $\mathbb{R}$.
    \begin{enumerate}
        \item 
        Let $s \in \mathbb{R}$. Show that $\displaystyle{\limsup_{n \to \infty} s_n = s}$ if and only if the following two conditions hold:
        \begin{itemize}
            \item 
            For all $\epsilon > 0$, there is $N \in \mathbb{N}$ so that $s_n < s + \epsilon$ for all $n \geq N$.

            \item 
            For all $\epsilon > 0$ and for all $N \in \mathbb{N}$, there exists $n \geq N$ so that $s_n > s - \epsilon$.\\
        \end{itemize}
        
        \item 
        Show that some subsequence of $\{s_n\}$ converges (in the sense of extended real numbers, if necessary) to $\displaystyle{\limsup_{n \to \infty} s_n}$ and some subsequence of $\{s_n\}$ converges to $\displaystyle{\liminf_{n \to \infty} s_n}$.\\

        \item 
        Define 
        \begin{equation*}
            E = \{x \in [-\infty, \infty] : \text{ some subsequence } \{s_{n_k}\} \text{ of } \{s_n\} \text{ converges to } x\}.
        \end{equation*}
        Show that $\displaystyle{\limsup_{n \to \infty} s_n = \sup E}$ and $\displaystyle{\liminf_{n \to \infty} s_n = \inf E}$.\\[5pt] \textit{Hint:} what happens to $\limsup$ and $\liminf$ on subsequences?\\
        
    \end{enumerate}

    \item 
    Let $a > 1$. In this problem, we will define what it means to write $a^r$ for a positive \textit{irrational} number $r$. We first define exponentiation for positive \textit{integer} exponents in the usual way, via repeated multiplication. We know this satisfies some properties, for example
    \begin{alignat*}{2}
        a^{r + s} &= a^r a^s \qquad \qquad \qquad &&(\ast)\\[5pt]
        (ab)^{r} &= a^r b^r \qquad \qquad \qquad &&(\Delta)
    \end{alignat*}
    for $a, b > 0$ and $r, s > 0$ integers. Now, we proved in class that for any positive integer $q$, there is a \textit{unique} positive number whose $q$th power is $a$, which we denote $\sqrt[q]{a}$. For a positive rational number $r = \frac{p}{q}$, we define $a^r = \sqrt[q]{a^p}$.
    \begin{enumerate}
        \item 
        Show that this definition makes sense, i.e. is independent of the representation of $r$ as a fraction. That is, show that if rational numbers $\frac{p}{q} = \frac{r}{s}$ then $\sqrt[q]{a^p} = \sqrt[s]{a^r}$. \textit{Hint:} use lowest terms and the uniqueness of $q$th roots.\\

        \item 
        Show that properties $(\ast)$ and $(\Delta)$ above still hold for positive rational numbers. \textit{Hint:} find a common denominator and use uniqueness again.\\

        \item 
        For positive rational $r < s < t$, show that $a^r < a^s$ and that
        \begin{equation*}
            0 < a^s - a^r = a^r(a^{s - r} - 1).
        \end{equation*}
        Conclude that for any positive rational $r, s$ and any rational $t$ larger than both $r$ and $s$, we have
        \begin{equation*}
            |a^r - a^s| < a^t(a^{|r - s|} - 1).
        \end{equation*}

        \item 
        Now let $r$ be a positive \textit{irrational} number. We know there is a sequence of rational numbers $q_k$ that converge to $r$ (for example, the decimal expansion of $r$). We define
        \begin{equation*}
            a^r = \lim_{k \to \infty} a^{q_k}.
        \end{equation*}
        Show that the sequence $\{a^{q_k}\}$ is Cauchy, and conclude that this limit exists.\\ \textit{Hint:} Use part (c) and Rudin Theorem 3.20(d).\\
        
        \item
        Show that the choice of sequence $\{q_k\}$ is irrelevant; that is, show that if $\{p_k\}$ and $\{q_k\}$ are sequences of rational numbers converging to $r$, then 
        \begin{equation*}
            \lim_{k \to \infty} a^{p_k} = \lim_{k \to \infty} a^{q_k}.\\
        \end{equation*}

        \item 
        Using limit laws, prove that properties $(\ast)$ and $(\Delta)$ continue to hold for positive irrational exponents.\\
        
    \end{enumerate}

    \textit{Remark:} For $r < 0$ we then define $a^r = \frac{1}{a^{-r}}$, and for $0 < a < 1$ we define $a^r = \left(\frac{1}{a}\right)^{-r}$. It is then a matter of checking cases to verify that all the properties of exponentiation we want are indeed true.
    
    
    
\end{enumerate}



\end{document}     